\subsection{First Letter}

\begin{quotex}
We have 11 letters from \textbf{Rene Guenon} to \textbf{Julius Evola} in the span from 1930 to 1951. In earlier letters to \textbf{Guido de Giorgio}, we saw Guenon's frustration with some of Evola's views. Here, Guenon confronts Evola directly, although more politely than in his letters to de Giorgio. Evola claims to have read all of Guenon's books; nevertheless, it seems he either misunderstood or rejected some significant aspects of Guenon's view. I believe there is some of both.

For example, Evola does not seem to understand Guenon's notion of the Intellect and thus shows little interest in ideas like the Supreme Identity. On the other hand, Evola is committed to his own philosophical system which differs from Guenon's metaphysical writings in important ways. In a letter to Mircea Elide, as we will soon see, Evola points out indeed that his doctrine is contained in his philosophy of the Absolute Individual.

The first of the letters follows:

\end{quotex}
24 August 1930

Cairo, Egypt

You must have thought that I would not respond to your letter, which reached me in Paris a little more than a year ago.

The truth is that at that time I was quite ill and, subsequently, different unforeseen difficulties and commitments of every type made me always delay every correspondence that was not absolutely urgent. Time flew by quickly and I never succeeded in doing everything I wanted to. I'm taking advantage of the fact that I am close to a little office in this residence to finally write you, asking you to excuse this excessive delay.

I have to tell you how little I was able to understand at all the interest that you showed in the reading of my books.

Obviously, the point of view you are assuming is quite distinctive and certainly cannot be mine, but I am pleased to see that that has not prevented you from getting rid of the anti-Oriental prejudice that, by your own admission, you used to hold. I wish that many others in the West would have the same attitude and come to understand the ancient doctrines of the Orient.

You ask me about [Jacques] Maritain; notwithstanding everything, I have always had friendly relations with him; as to ideas, we are in agreement especially on a negative point of view, that is, on ``anti-modernity". Apart from that, even he, disgracefully, is full of prejudices against the Orient; at least he was, because it seems that those prejudices since a short time ago have been attenuated; but, something strange, it is fed by a type of fear in the face of what one does not know, and it is a disagreeable thing, because it prevents him from broadening his own point of view.

But permit me to point out to you, from the moment that you read all my books that, after The Crisis of the Modern World, there is another, Spiritual Authority and Temporal Power, what was published last year.

Currently, I am working on The Symbolism of the Cross that will definitely be published toward the end of this year.

Excuse the briefness of my letter; I would like to be able to more or less get up to date with our correspondence.



\flrightit{Posted on 2012-07-16 by Cologero }

\begin{center}* * *\end{center}

\begin{footnotesize}\begin{sffamily}



\texttt{Avery on 2012-07-17 at 00:38 said: }

It is surprising that Evola would have once held a prejudice against the East, since his work ends up being much more relevant to East Asian traditions than Guenon's.


\hfill

\texttt{Cologero on 2012-07-17 at 07:12 said: }

It's true that Evola adopted aspects of Tantrism, Buddhism, and Taoism into his own system but, unlike Guenon, Evola's integration of those Eastern doctrines was idiosyncratic. I suspect Guenon was referring to Evola's Euro-centrism and white supremacism; after all, Evola continued to support European colonialism.


\hfill

\texttt{Umisumi on 2015-12-20 at 09:47 said: }

Good work of translation. May i ask you where you obtained these letters? Do you have copies or scans of the originals?


\hfill

\texttt{Cologero on 2015-12-20 at 10:09 said: }

\url{http://www.libreriaeuropa.it/scheda.asp?id=5216\&ricpag=1}

Evola saved all his correspondence, unlike Guenon, so we just have one half of the correspondence. That was the era before photocopy machines. The letters were handwritten, apparently, so Evola didn't keep carbon copies of his side (if anyone still remembers those). Evola wrote in Italian; Guenon responded in French.

An interesting point that people today tend to overlook is that even photos were had to come by. After years of correspondence and even Evola's translations of Guenon into Italian, Evola had no idea what Guenon looked like.


\end{sffamily}\end{footnotesize}
